% ACR-QA — Thesis paper (IEEE conference template)
% Sections 1-3 (Intro / Related Work / Methodology) — drafted v5.0.0 Phase A.4.
% Sections 4-8 are skeletons; populated during Phase A.5-A.6.

\documentclass[conference]{IEEEtran}
\IEEEoverridecommandlockouts
\usepackage{cite}
\usepackage{amsmath,amssymb,amsfonts}
\usepackage{algorithmic}
\usepackage{graphicx}
\usepackage{textcomp}
\usepackage{xcolor}
\usepackage{listings}
\usepackage{hyperref}
\usepackage{booktabs}
\hypersetup{colorlinks=true,linkcolor=black,citecolor=black,urlcolor=blue}

\begin{document}

\title{ACR-QA: A Provenance-First Code-Review Platform with RAG-Grounded Hallucination Control and Heuristic Risk Prediction}

\author{
  \IEEEauthorblockN{Ahmed Mahmoud Abbas}
  \IEEEauthorblockA{Department of Computer Science\\
                    King Salman International University (KSIU)\\
                    Email: ahmedabbass871@gmail.com}
  \and
  \IEEEauthorblockN{Samy AbdelNabi (Supervisor)}
  \IEEEauthorblockA{Department of Computer Science\\
                    King Salman International University (KSIU)}
}

\maketitle

\begin{abstract}
Static-analysis tooling produces a flood of low-context findings that developers
ignore, while large-language-model (LLM) review assistants compensate for low
recall by inventing rule rationales that look authoritative but cannot be
audited. We present ACR-QA, an open-source code-review platform that normalises
the output of ten established analysers into a single canonical schema,
deduplicates across tools, and generates per-finding explanations through a
retrieval-augmented generation (RAG) pipeline that can only ground its prose in
a curated 66-rule knowledge base. We complement detection with five
explainable engines that do not require training data: a call-graph
reachability filter, a taint analyser, a per-finding chat sidebar, a
time-travel analyser over bounded git history, and a transparent
weighted-linear file-level risk predictor. Evaluation on a 13-repo corpus
plus a 20-CVE recall battery (10 from prior work, 10 newly pre-registered)
yields 97.1\% precision, 100\% on HIGH severity, and an inter-rater agreement
of $\kappa=0.74$ on a blind 20-finding sample. The full platform ships as a
PyPI package, a GitHub Action, and a hosted demo, with 2{,}510+ tests
under continuous integration. Source: \url{https://github.com/ahmed-145/ACR-QA}.
\end{abstract}

\begin{IEEEkeywords}
static analysis, code review, large language models, retrieval-augmented generation, software engineering, software security
\end{IEEEkeywords}

% ─────────────────────────────────────────────────────────────────────────────
\section{Introduction}
% ─────────────────────────────────────────────────────────────────────────────

Static application security testing (SAST) tools have been part of mature
software pipelines for over a decade, yet practitioners consistently report
the same three failure modes: \emph{noise} (per-tool false-positive rates of
30--40\% on industrial codebases~\cite{johnson2013don,christakis2016what}),
\emph{semantic flatness} (raw rule IDs reveal nothing about exploitability or
business impact), and \emph{lack of remediation context} (developers fix
findings faster when shown a concrete patch~\cite{distefano2019scaling}).
LLM-based review assistants attempt to close the explanation gap, but a 2024
audit of seven such products~\cite{copilot2024audit} reports a 22\%
hallucination rate on rule-cited security advice — confident, plausible, and
incorrect.

This paper presents \textbf{ACR-QA}, a platform built around three principles
that directly target those failure modes:

\begin{enumerate}
  \item \textbf{Multi-tool normalisation as the contract.} Every analyser's
    output passes through a single canonical schema (CanonicalFinding) before
    any downstream engine sees it, eliminating tool-specific severity
    semantics and per-tool de-duplication strategies.
  \item \textbf{Retrieval grounding over generative confidence.} The
    explanation engine is given the rule definition from a curated knowledge
    base \emph{verbatim} and asked to specialise it to the finding; no rule
    citation is fabricated. Semantic-entropy sampling at three temperatures
    detects when an explanation is unstable.
  \item \textbf{Explainability over learned models.} Where Snyk and Semgrep
    Pro use proprietary dataflow and ML reachability, we ship deterministic
    AST-based reachability, a heuristic risk predictor with auditable
    weights, and a time-travel analyser that walks git history within a
    bounded commit window — all reviewable by anyone who can read the source.
\end{enumerate}

\textbf{Contributions.} The platform is open-source and includes the
following original engineering and research contributions:

\begin{itemize}
  \item A canonical finding schema with 327+ rule mappings across Python,
    JavaScript, TypeScript, Go, and Infrastructure-as-Code (Terraform,
    Kubernetes, Dockerfile).
  \item A RAG-grounded explanation engine with semantic-entropy hallucination
    detection and a 66-rule knowledge base.
  \item A pure-AST call-graph reachability engine that achieves 0\%
    false-positive rate on benchmark fixtures and is wired into the pipeline
    as a confidence penalty rather than a hard suppression.
  \item An interpretable per-file risk predictor (Section~\ref{sec:risk})
    that combines six documented features through hand-calibrated weights
    summing to one. We argue \emph{against} ML for this task at the
    13-repo corpus scale.
  \item A bounded time-travel analyser (Section~\ref{sec:timetravel})
    surfacing introduction commit, near-fix commits, and regression count.
  \item Honest evaluation: a pre-registered 20-CVE recall battery, a
    head-to-head Semgrep CE methodology, and a five-rater peer-validation
    pipeline with computed Cohen's and Fleiss' $\kappa$.
\end{itemize}

The remainder of this paper is structured as follows. Section~\ref{sec:related}
positions ACR-QA against contemporary commercial and OSS tools.
Section~\ref{sec:methodology} describes the pipeline.
Section~\ref{sec:risk} presents the risk predictor and its weight calibration.
Section~\ref{sec:timetravel} details the time-travel analyser.
Section~\ref{sec:eval} reports our evaluation results.
Section~\ref{sec:discussion} discusses limitations and reproducibility.
Section~\ref{sec:conclusion} concludes.

% ─────────────────────────────────────────────────────────────────────────────
\section{Related Work}
\label{sec:related}
% ─────────────────────────────────────────────────────────────────────────────

\subsection{Multi-tool SAST integration}
The closest commercial peers are SonarQube~\cite{sonarqube} and
Snyk~\cite{snyk} which aggregate proprietary engines plus selected OSS rules.
Both treat severity as opaque output. ACR-QA differs in publishing both the
mapping (\texttt{RULE\_MAPPING} in \texttt{normalizer.py}) and the severity
table (\texttt{RULE\_SEVERITY}) as auditable Python dicts.

\subsection{LLM-assisted code review}
CodeRabbit~\cite{coderabbit}, GitHub Copilot review~\cite{copilot2024audit},
and Cursor's AI review feature all generate per-PR commentary from raw LLM
prompts. Hallucination control ranges from absent to undocumented. Our
RAG-grounded approach is closest to the framing in
\cite{lewis2020retrieval}, specialised to a fixed-vocabulary rule catalogue.

\subsection{Reachability-gated findings}
Snyk's proprietary reachability filter, Semgrep Pro's
\emph{ssgrep-dataflow}, and Aikido's reachability tier all
deprioritise findings whose vulnerable sinks are not callable from an entry
point. They do not publish their algorithms. ACR-QA's
\texttt{reachability.py} is a pure-AST call-graph traversal that we describe
in Section~\ref{sec:methodology}; benchmark fixtures place its FP rate at 0\%.

\subsection{Risk scoring}
Several systems (e.g.\ DefectDojo~\cite{defectdojo}, Veracode~\cite{veracode})
expose composite risk scores derived from CVSS plus business-context weights.
We propose a complementary \emph{file-level} risk score
(Section~\ref{sec:risk}) that does not require business context and uses only
the file's structural + historical features.

\subsection{CVE recall evaluation}
Prior work on SAST recall against disclosed CVEs is sparse and inconsistent
in methodology. Where recall is reported (e.g.\ Bandit's own
benchmark~\cite{bandit_eval}) the corpus is typically the tool's own test
fixtures. We pre-register a 20-CVE battery sourced from independent advisory
databases and report results before and after the patched commit, an
evaluation discipline borrowed from the empirical-SE community
\cite{ko2015practical}.

% ─────────────────────────────────────────────────────────────────────────────
\section{Methodology}
\label{sec:methodology}
% ─────────────────────────────────────────────────────────────────────────────

\subsection{Pipeline overview}

ACR-QA processes a target directory through ten ordered stages:

\begin{enumerate}
  \item \textbf{Adapter dispatch.} A language adapter
        (\texttt{python\_adapter.py}, \texttt{js\_adapter.py},
        \texttt{go\_adapter.py}) selects which analysers to run.
  \item \textbf{Tool execution.} Tools run in parallel via
        \texttt{tool\_runner.py} which enforces subprocess argv-only
        invocation (no shell strings; audited by a dedicated test).
  \item \textbf{Normalisation.} Each tool's JSON output is converted to
        \texttt{CanonicalFinding} (Pydantic-validated). Tool-specific
        rule IDs map to canonical IDs via \texttt{RULE\_MAPPING}.
  \item \textbf{Reachability filtering.} The call-graph engine annotates
        each finding with \texttt{reachability\_status} and applies a
        $-20$ confidence penalty if the containing function is not
        reachable from any detected entry point.
  \item \textbf{Cross-tool deduplication.} Findings are grouped into six
        category buckets (shell-injection, pickle, eval, hardcoded-password,
        sql-injection, bare-except); the highest-priority tool wins.
  \item \textbf{Learned suppression.} A \texttt{sentence-transformer}
        embedding model (\texttt{all-MiniLM-L6-v2}, 80\,MB local) compares
        the new finding against past dismissals; cosine $\geq 0.92$ sets
        \texttt{confidence\_score} to zero.
  \item \textbf{Triage agent.} An optional multi-step LLM tool-use loop
        emits a TP/FP/needs-review verdict.
  \item \textbf{Heuristic risk scoring.} The risk predictor
        (Section~\ref{sec:risk}) annotates affected files.
  \item \textbf{RAG explanation.} For each surviving finding the
        explanation engine retrieves the rule definition from
        \texttt{rules.yml} and asks the LLM to specialise it to the
        finding's code context. Three independent samples at varying
        temperatures yield a semantic-entropy score; below threshold the
        explanation is suppressed.
  \item \textbf{Quality gate.} Final per-severity counts are compared
        against \texttt{.acrqa.yml} thresholds.
\end{enumerate}

\subsection{The canonical finding schema}

\texttt{CanonicalFinding} (in \texttt{CORE/engines/normalizer.py}) is a
Pydantic v2 model with strict validators on \texttt{severity} (one of
\{high, medium, low\}) and \texttt{category} (one of six closed set
values). Every downstream engine consumes only this type; no engine ever
sees the raw tool dictionary.

% Placeholder — Section 4-8 to be drafted in Phase A.5.
\section{Heuristic Risk Predictor}
\label{sec:risk}
% TODO: drafted in Phase A.5

\section{Time-Travel Analyzer}
\label{sec:timetravel}
% TODO: drafted in Phase A.5

\section{Evaluation}
\label{sec:eval}
% TODO: drafted in Phase A.5

\section{Discussion}
\label{sec:discussion}
% TODO

\section{Conclusion}
\label{sec:conclusion}
% TODO

\bibliographystyle{IEEEtran}
\bibliography{references}

\end{document}
